\section{Mass Is The Strain}

Charge counts the completed loop. Mass asks a different question. It asks
what the loop costs when the reading is pushed past threshold. A count can
tell the grammar that traversal has returned. It cannot by itself tell the
grammar how the return resists variation. Resistance appears only when the
grammar reads not merely a difference, but a difference of differences.

The first trip is null. The grammar tests a reading and finds no admitted
motion. The level stays where it was. This does not mean nothing has been
learned. The null trip establishes a baseline: under this much pressure, the
record permits no displacement that the reader can carry. A baseline is a
reading, not an absence. It tells the grammar what counts as no motion under
the current trial.

The second trip reaches threshold. Now a slip begins. The reading has crossed
from not moving to moving, from pooled to selected, from held to displaced.
Threshold is not yet mass. It is the first place where the old balance no
longer absorbs the variation. The grammar reads that a transition has become
admissible, but it has not yet read how the transition itself changes under
the next push.

The third trip supplies the response. After null and threshold, the grammar
can compare the change in change. It can ask how the displacement differs
from the threshold displacement, how the response bends rather than merely
starts, how the carried loop resists further variation. This is the second
difference. It is the first reading rich enough to carry strain.

Strain is the grammar's name for that carried bending. A first difference
reads displacement. A second difference reads the variation of displacement.
If the second difference is zero, the response has no readable bend at this
scale. If it is nonzero, the loop has carried a cost that cannot be erased by
renaming the path. The grammar reads that cost as mass.

This reverses a common comfort. The comfortable story begins with mass as a
thing and then explains how it resists acceleration. The grammar begins with
a finite response and asks what must be carried if the response bends past
threshold. Mass is not injected into the ledger as a hidden weight. It
surfaces when the reading reaches the strain-bearing face of the loop. The
name mass belongs to the second-difference strain, not to a prior substance
waiting behind it.

The relation to inertia is therefore grammatical before it is physical. A
reading that preserves its state under the null trip, begins to move at
threshold, and resists the next change has acquired an inertial face. The
resistance is not separate from the strain. It is the strain read as response
to variation. In this disciplined sense, inertial response and gravitational
strain can be treated as two readings of one carried cost. The chapter need
not claim a complete physical theory of gravity to say that the grammar
identifies response and strain at the face where the second difference is
read.

The fluxion example illustrates the finite source of the reading. A smooth
slope can be treated as though it came from a vanished increment, but the
grammar does not need that ghost. It reads finite anchor differences and
then forces the smooth shadow that preserves them under refinement. The
second difference belongs to the same discipline. It is not an infinitesimal
object. It is the stable reading forced by finite comparisons of response.

The sombrero example illustrates threshold and remembered orientation. In a
symmetric field, many orientations may seem equally available. Once a
particular minimum is selected, the record must remember which orientation
was chosen. The cost of maintaining that remembered choice appears as
curvature around the selected state. The example is only a lantern. The
chapter's claim is the simpler one: after threshold, a selected response
that bends carries strain, and that strain is the mass face.

The three-trip grammar also explains why charge and mass belong to one
number without collapsing into one face. Charge counts how the loop comes
home. Mass reads how the loop strains when the count is varied past
threshold. Charge is a lower loop tally; mass is an upper response reading.
The same closed loop carries both, but it carries them under different
questions. A charge without strain is an orientation count. A mass without
charge would be a bending response without the loop count that places it in
the ledger.

The backward walk needs both. If it returned only with charge, it would name
a traversal but not the resistance of the traversal. If it returned only with
strain, it would name a cost but not the closed count that made the cost
billable. The grammar therefore brackets the number. The lower face reads
charge. The upper face reads strain. The second difference past threshold is
where mass enters the bracket.

With charge and mass in the bracket, the next question is forced. A number
that carries loop count and strain is still not exhausted. It can also be
read by phase, because orientation can be transported and decided, and it
can be read by sameness, because closure touches the needle. The one number
therefore has four faces.
